\documentclass[10pt, twocolumn, landscape, a4paper]{article}

\usepackage[margin=1.5cm]{geometry}
\usepackage[utf8]{inputenc}
\usepackage[spanish,es-noshorthands]{babel}
\usepackage{amssymb, amsmath, amsbsy}
\usepackage{verbatim}
\usepackage{mathpazo}
\usepackage{tasks}
\usepackage{enumerate}
\usepackage{enumitem}
\usepackage{graphicx}
\usepackage{tikz} 
% \usetikzlibrary{angles,quotes}
\usepackage[pdftex]{hyperref}
\hypersetup{pdftitle={Matriz de transformación bidimensional (2D)},pdfauthor={Luis Carrillo Gutiérrez}}

\fontfamily{cmss}\selectfont 
\renewcommand{\familydefault}{\sfdefault}
\pagestyle{empty}
\setlength{\columnsep}{1.2in}

\begin{document}
\section*{Matriz de transformación bidimensional (2D)}
Sea un objeto (un cuadrilátero) de 4 vértices/puntos ($P_{n}$) :
\begin{eqnarray*}
P_{\,1} &=& (-1, \;1) \\
P_{\,2} &=& (\;1, \;1) \\
P_{\,3} &=& (\;1, -1) \\
P_{\,4} &=& (-1, -1) 
\end{eqnarray*}

Gráfico: 
\begin{center}
\begin{tikzpicture}
	\draw[help lines,step=1cm] (-2,-2) grid (2,2);
	
	\coordinate (p1) at (-1, 1); 
	\coordinate (p2) at (1, 1); 
	\coordinate (p3) at (1, -1); 
	\coordinate (p4) at (-1, -1); 
	
	\draw[line width=1.75] (p1) -- (p2) -- (p3) -- (p4) -- cycle;
	\draw (p1) circle (.075);
	\draw (p2) circle (.075);
	\draw (p3) circle (.075);
	\draw (p4) circle (.075);
	
	\draw (-1.35, 1.3) node {$P_{\,1}$};
	\draw (1.35, 1.3) node {$P_{\,2}$};
	\draw (1.35, -1.3) node {$P_{\,3}$};
	\draw (-1.35, -1.3) node {$P_{\,4}$};
	% \draw (.25, .25) node {$(0,\,0)$};
\end{tikzpicture}
\end{center}

\noindent Un vértice/punto bidimensional se representa mediante sus coordenadas (par ordenado) con un orden estricto $(x,\,y)$, es decir como un \textbf{vector} que indica el desplazamiento horizontal y luego el vertical. Además se puede representar en su \textbf{forma homogénea}, como una matriz y en concreto como una {\tt matriz columna} o {\tt vector columna} (de orden $n \times 1$). Entonces :

$$P_{\,1} = (x,\,y) = \begin{bmatrix}x \\ y \end{bmatrix}_{2 \times 1} \implies (-1,\,1) = \begin{bmatrix} -1 \\ 1 \end{bmatrix}_{2 \times 1}$$
$$P_{\,2} = (1,\,1) = \begin{bmatrix} \,1 \\ \,1 \end{bmatrix}$$
$$P_{\,3} = (1,\,-1) = \begin{bmatrix} 1 \\ -1 \end{bmatrix}$$
$$P_{\,4} = (-1,\,1) = \begin{bmatrix} -1 \\ 1 \end{bmatrix}$$

\noindent Esto no es más que la transpuesta de un \textbf{vector fila} $(\,x\;\;y\,) = [\,x\;\;y\,]$, por tanto nos referimos a una coordenada como también la transpuesta del vector fila : $[\,x\;\;y\;\;z\;\;w\,]^{T}$.

\subparagraph*{¿Qué son las transformaciones geómetricas?} --

\noindent Son cambios de posición, tamaño u orientación sobre objetos vectoriales (en concreto sobre las coordenadas de sus vértices). Estos cambios/transformaciones incluyen escalar, rotar, trasladar, sesgar (deformar), o alguna combinación de todas las mencionadas.

\subparagraph*{¿Cómo se transforman objetos en el espacio bidimensional?} --

\noindent Mediante el uso de matrices (en concreto una \emph{matriz de transformación}), que es el mecanismo matemático para entender cómo transformar geométricamente un objeto. 

\subparagraph*{¿Por qué usar matrices?} --

\noindent Una matriz puede ser usada para describir o calcular transformaciones en un espacio bidimensional (2D) de forma general. Su uso se promueve para describir cualquier /transformación afín/ a un objeto expresado en un grupo de vértices (o coordenadas). 

\subparagraph*{¿Qué es una matriz de transformación?} --

\noindent Es una matriz cuadrada, no menor a una matriz de orden $2\times\!2$, que se usa para transformar/modificar la posición, orientación, forma o tamaño de un objeto (mediante operaciones entre matrices /multiplicación/), aprovechando que las coordenadas de un objeto se puede expresar como una matriz columna. \\

\noindent Por tanto las \emph{transformaciones bidimensionales} pueden ser representadas/expresadas por una matriz de orden $2 \times 2$; sin embargo, la traslación no encaja bien dentro de está regla, por lo tanto una matriz de orden $3 \times 3$ es requerida, y por ello es que las coordenadas de los vértices de los objetos se expresen usando {\tt coordenadas homogéneas} bidimensionales o una matriz de orden $3\times\!1$. 

$$P_{\,1} = (x,\,y) = \begin{bmatrix}x \\ y \\ \mathtt{1}\end{bmatrix}_{3 \times 1} \implies (-1,\,1) = \begin{bmatrix} -1 \\ 1 \\ 1\end{bmatrix}_{3 \times 1}$$

% al escalarlo
\noindent Es lo que llamamos {\tt normalizar} el vector o reducir a la unidad la tercera fila. Por ejemplo, convertir $[\,4\;\;6\;\;2\,]^T$ a $[\,2\;\;3\;\;1\,]^T$. La excepción es un punto al infinito (un vector direccional), en cuyo caso la tercera fila es 0. \\ 

\noindent Si al multiplicar la matriz de transformación por un vector columna/matriz columna normalizado, se obtiene otro vector columna normalizado (transformado),  siendo esta la forma de representar la transformación de puntos en 2D mediante matrices de orden 3×3 :

$$
\begin{vmatrix} x\,' \\ y\,' \\ z\,' \end{vmatrix}_{3x1} = \begin{bmatrix} a & b & c \\ d & e & f \\ g & h & i \end{bmatrix}_{3x2} \times \begin{bmatrix} x \\ y \\ z \end{bmatrix}_{3x1} 
$$
\end{document}